% Four-unit course map -- SHARED SOURCE.
%
% \input by two places; do not duplicate the picture:
%   ../slides/CHEME-5660-L1a-Slides-Fall-2026.tex   (typeset into the deck)
%   Fig-L1a-CourseMap.tex                           (standalone -> PDF -> SVG)
% Editing this file updates both. Rebuild with `make` in this directory.
%
% DESIGN, following templates/vnflow.sty and the L1b abstract-asset timeline:
%
%   SPINE.   A real week axis carries the figure, exactly as the L1b timeline
%            hangs its cash flows off a time axis. The units are not floating
%            boxes joined by arrows; they sit ON the semester.
%   COLOUR.  Semantic, not decorative. vnlink is reserved for the CAPABILITY
%            each unit adds -- the payload of the figure. Structure (axis,
%            boxes, ticks) is vnink/vnmuted, and nothing else is coloured.
%   GEOMETRY. Every box uses the same fixed-height title area, and every
%            capability line uses the same text width, so the four columns
%            align on a shared grid. No label alternates above and below; each
%            sits at a fixed offset from its own box.
%   PAYLOAD. The brace states the question the whole course answers, so the
%            figure teaches "these four together answer one thing" rather than
%            "first this, then that."
%
% The card titles below are deliberately shorter than the workbook unit names:
% this first-day map names the subject of each unit in plain language. The
% capability line underneath says what students will be able to do with it.
%
% Unit 1 uses "Treasury securities" rather than the broader "payments over
% time" because Treasury securities are the only asset class studied in that
% unit. This deliberately introduces the name before the five-promises reveal.
\begingroup
\hyphenpenalty=10000 \exhyphenpenalty=10000 \pretolerance=10000
\begin{tikzpicture}[
    font=\sffamily\small,
    >=Latex,
    unit/.style     = {draw=vnink,line width=0.9pt,rounded corners=2pt,
                       fill=vnsoft,minimum width=3.45cm,minimum height=1.5cm,
                       align=center,text width=3.15cm,inner sep=3pt},
    adds/.style     = {text=vnlink,font=\sffamily\small,align=center,
                       text width=3.45cm},
    weeks/.style    = {text=vnmuted,font=\sffamily\footnotesize},
    tick/.style     = {vnmuted,line width=0.7pt}]

  % ---- the four units, on one grid ---------------------------------
  \foreach \i/\x/\utitle in {%
      1/0/{Treasury securities},
      2/4.35/{Stocks and portfolios},
      3/8.70/{Options and other derivatives},
      4/13.05/{Financial decision-making}}{%
    \node[unit] (u\i) at (\x,0)
      {\textcolor{vnink}{\bfseries Unit \i}\\[5pt]
       \parbox[c][0.62cm][c]{3.05cm}{\centering\footnotesize \utitle}};}

  % ---- the spine: the semester ------------------------------------
  \draw[vn timeline] (-1.95,-1.35) -- (15.6,-1.35) node[right,text=vnink] {weeks};
  \foreach \i/\x/\wk in {1/0/{1--2}, 2/4.35/{3--7}, 3/8.70/{8--11}, 4/13.05/{12--15}}{%
    \draw[tick] (\x,-1.25) -- (\x,-1.45);
    \node[weeks,below=1pt] at (\x,-1.45) {weeks \wk};}

  % ---- what each unit adds: one aligned band, one shared baseline --
  \foreach \i/\x/\cap in {%
      1/0/{compare payments\\across dates},
      2/4.35/{build portfolios\\under uncertainty},
      3/8.70/{value investments\\tied to other values},
      4/13.05/{build and test\\financial strategies}}{%
    \node[adds,anchor=north] at (\x,-2.15) {\cap};}

  % ---- the payload: one question, four units ----------------------
  \draw[vn brace] (-1.78,1.05) -- (14.83,1.05)
    node[midway,above=7pt,text=vnink]
    {\bfseries Given investment choices, uncertainty, and your priorities: what should you do?};
\end{tikzpicture}
\endgroup
